\subsection{What is needed for a stronger Mertens bound}\label{disc:mertens-improvement}

The issue is not another representation of $M(n)$, but an estimate for one of its representations that uses additional cancellation. The classical comparison in this paper is \eqref{arith:mertens}: with
\[
 V(n)=\frac{(\log n)^{3/5}}{(\log\log n)^{1/5}},
\]
it gives $|M(n)|\le Cn\exp(-cV(n))$ for suitable positive constants and all sufficiently large $n$. A new proof of $M(n)=o(n)$, or of a weaker quantitative bound, would not improve this estimate. An explicit improvement must specify its constants and range; a stronger asymptotic scale must specify its rate. For example, proving $|M(n)|\le Cn\exp(-a(n))$ with $a(n)/V(n)\to\infty$ would improve the displayed scale. No such result is established here.

There is a concrete sufficient condition that does not require computing a determinant. Let $\varepsilon(n)>0$ be a proposed error rate. Construct a vector $z_n\in\mathbb R^{n-1}$ and bound its residual by
\begin{equation}\label{disc:certificate}
 \norm{\one_n-H_n^Tz_n}_2^2
 \le \frac{n^2\varepsilon(n)^2}{Q(n)}.
\end{equation}
Then $|M(n)|\le n\varepsilon(n)$. Indeed, $H_n\muvec_n=0$ implies
\[
 M(n)=\muvec_n^T(\one_n-H_n^Tz_n),
\]
and Cauchy--Schwarz gives the conclusion. Thus the missing ingredient could be an explicit choice of row coefficients together with a uniform residual estimate. Choosing the exact least-squares solution and using \eqref{geo:height} to evaluate its error merely recovers $M(n)^2/Q(n)$; it supplies no independent cancellation estimate. Since $Q(n)\sim c_0n$, the required squared residual has scale $n\varepsilon(n)^2/c_0$, with constants controlled if an explicit bound is intended.

Two other sufficient certificates are already available from the geometric identities. For a choice $y=y(n)\ge2$, it suffices to prove
\begin{equation}\label{disc:component-certificate}
 F(n,y)\le\frac{n^2\varepsilon(n)^2}{Q(n)}.
\end{equation}
Alternatively, for a choice $\tau=\tau(n)>0$, it suffices to prove
\begin{equation}\label{disc:resolvent-certificate}
 \tau\one_n^T(H_n^TH_n+\tau I_n)^{-1}\one_n
 \le\frac{n^2\varepsilon(n)^2}{Q(n)}.
\end{equation}
The first follows from \eqref{geo:variance}, the second from \eqref{geo:resolvent}. Each is a target for further work, not an estimate proved by writing the formula. In particular, the component variance identity shows that $F$ includes both the global Mertens contribution and nonnegative variation between components. A useful upper bound must control both. The fixed-power profile has $F(n,n^{1/u})\sim c_0nJ(u)$ with $J(u)>0$ for every fixed $u\ge1$; that regime cannot certify $\varepsilon(n)\to0$. Obtaining an explicit decay rate as $u$ grows requires additional quantitative control; the qualitative cutoff in Theorem~\ref{geo:profile} follows from the fixed-parameter limits and monotonicity. The displayed resolvent depends on the weights of $\one_n$ in the eigenspaces; the eigenvalue information established here alone does not control its smallness.

The singular-value approach has the same unresolved step. The formula for the exceptional smallest singular value already contains $|M(n)|$. Combining that formula with the product of the other singular values therefore reproduces $|M(n)|$, rather than bounding it independently. To gain arithmetic information, one needs an independently obtained estimate for the exceptional direction, or uniform information on a growing part of the spectrum that controls the full product without importing the desired Mertens bound. Fixed-index limits, the sum of squared singular values, and the product with the smallest factor omitted do not provide that missing estimate.

Any proposed improvement should therefore state three things separately: the new arithmetic or geometric input; the uniform argument transferring it to one of these certificates; and the resulting rate for $M(n)$ compared with \eqref{arith:mertens}. Bounds obtained by inserting the existing Mertens estimate into restricted sums are useful consequences, but cannot themselves serve as independent evidence for stronger cancellation.
